\section*{Hoofdstuk \ref{ch:b1:chemical-biogenic-sediments} --- Carbonaat-, kiezel- en verdampingssedimenten}

\begin{solution}{exo:b1:chemical-biogenic-sediments:1}
Chloride $0.5503\times35.165 = \qty{19.35}{g}$; natrium $0.3066\times35.165 =
\qty{10.78}{g}$; sulfaat \qty{2.711}{g}; magnesium \qty{1.284}{g}. Natrium en
chloride samen maken $0.5503 + 0.3066 = 0.857$, dat is \qty{85.7}{\%} van het
zout.
\end{solution}

\begin{solution}{exo:b1:chemical-biogenic-sediments:2}
De fractie die overblijft is $1/\mathrm{CF}$: calciumcarbonaat $1/1.8 = 0.556$ blijft,
\qty{44.4}{\%} verdampt; gips $0.263$ blijft, \qty{73.7}{\%} verdampt; haliet $0.094$ blijft,
\qty{90.6}{\%} verdampt; magnesiumsulfaat $0.014$ blijft, \qty{98.6}{\%} verdampt.
\end{solution}

\begin{solution}{exo:b1:chemical-biogenic-sediments:3}
(a)~$\rho_s(1-\phi) = G/a = 8.4/0.0042 = \qty{2000}{kg.m^{-3}}$. (b)~Met
aragoniet, $1 - 2000/2930 = 0.32$, een porositeit van \qty{32}{\%}; met calciet
$1 - 2000/2710 = 0.26$. (c)~$a = 3.7/2000 = \qty{1.85e-3}{m}$, \qty{1.85}{mm} per
jaar.
\end{solution}

\begin{solution}{exo:b1:chemical-biogenic-sediments:4}
Tot \qty{4000}{m} daalt het gehalte slechts langzaam (88 naar 80); beneden
\qty{4000}{m} daalt het steil: het \omterm{def:b1:chemical-biogenic-sediments:ccd}{lysoklien} ligt op ongeveer \qty{4000}{m}. Het
gehalte zakt onder \qty{10}{\%} tussen \qty{4700}{m} (12) en \qty{4800}{m}
(6): $4700 + 100\times(12-10)/(12-6) = \qty{4733}{m}$, een CCD op ongeveer
\qty{4730}{m}, enkele \qty{700}{m} onder het \omterm{def:b1:chemical-biogenic-sediments:ccd}{lysoklien}.
\end{solution}

\begin{solution}{exo:b1:chemical-biogenic-sediments:5}
(a)~\ce{Ca^2+ + 2 HCO3- -> CaCO3 + CO2 + H2O}: één mol \ce{CO2} per mol \ce{CaCO3}, dus $44.01/100.09 = \qty{0.440}{kg}$ \ce{CO2} per kilogram \ce{CaCO3}. (b)~$3.7\times0.440 = \qty{1.63}{kg}$ \ce{CO2} per vierkante meter en per jaar. (c)~Van de twee gebruikte waterstofcarbonaationen gaat slechts één koolstofatoom in het \omterm{def:b1:rock-forming-minerals:mineral}{mineraal}; het andere verlaat als koolstofdioxide. Verkalking haalt alkaliniteit uit het water en verlaagt zijn vermogen om opgelost koolstof vast te houden; hoeveel van het \ce{CO2} de lucht bereikt, hangt af van de buffering van zeewater, een kwestie van de carbonaatevenwichten uit de scheikundeboeken.
\end{solution}

\begin{solution}{exo:b1:chemical-biogenic-sediments:6}
(a)~$t_c = (2000/350)^2 = \qty{32.7}{Myr}$; bij \qty{25}{km} per miljoen jaar,
$25\times32.7 = \qty{816}{km}$ van de as. (b)~$t_c = (2600/350)^2 =
\qty{55.2}{Myr}$, \qty{1380}{km}. (c)~Een vloer van \qty{60}{Myr} oud heeft rode
klei ontvangen gedurende $60 - 32.7 = \qty{27.3}{Myr}$ in het eerste geval, en slechts
$60 - 55.2 = \qty{4.8}{Myr}$ in het tweede.
\end{solution}

\begin{solution}{exo:b1:chemical-biogenic-sediments:7}
(a)~Ja: $a = 4.2 > r = 3$, en het \omterm{def:b1:chemical-biogenic-sediments:reef}{rif} wint \qty{1.2}{mm} per jaar op de zee
tot zijn top het oppervlak bereikt. (b)~$D$ groeit met $10 - 4.2 = \qty{5.8}{mm}$ per
jaar; van 2 tot \qty{50}{m} duurt het $48/0.0058 = \num{8300}$ jaar. (c)~Omdat
zijn groeisnelheid daalt naarmate het water dieper wordt: minder licht bereikt de algen die
in de koralen leven en hen voeden, dus daalt $a$ en raakt het \omterm{def:b1:chemical-biogenic-sediments:reef}{rif} steeds sneller achterop.
\end{solution}

\begin{solution}{exo:b1:chemical-biogenic-sediments:8}
(a)~De reactie is \ce{CaSO4.2H2O -> CaSO4 + 2 H2O}. (b)~De molaire volumes zijn $172.17/2.30 =
\qty{74.9}{cm^3}$ en $136.14/2.97 = \qty{45.8}{cm^3}$, zodat de vaste stof
$1 - 45.8/74.9 = \qty{38.8}{\%}$ van haar volume verliest. (c)~$2\times18.02/172.17 = 0.209$:
\qty{20.9}{\%} van de massa verlaat als water. Op diepte vrijgekomen, verhoogt dit water
de druk in de poriën van de omringende gesteenten; onder een afdichting kan het de gesteenten
zwakker maken en het glijden van breuken helpen, en het kan andere zouten oplossen en meevoeren.
\end{solution}

\begin{solution}{exo:b1:chemical-biogenic-sediments:9}
(a)~$9.2/255 = \qty{3.6}{\%}$ bewaard, \qty{96.4}{\%} opgelost. (b)~Waar de
productie zeer hoog is, is een paar procent ervan nog steeds een grote stroom, en waar
weinig klei of carbonaat het verdunt: onder de opwellende wateren van de Zuidelijke
Oceaan, de Noord-Pacifische Oceaan en de equatoriale Stille Oceaan, en onder de CCD.
(c)~Dunne, broze frustules lossen eerst op en robuuste blijven bewaard: het sediment
bewaart een selectie van de sterkste soorten, niet de gemeenschap die in het water leefde.
\end{solution}

\begin{solution}{exo:b1:chemical-biogenic-sediments:10}
(a)~$EA = 1.5\times\qty{1e10}{m^2} = \qty{1.5e10}{m^3}$ per jaar. (b)~Met
$Q_{\mathrm{in}} = \qty{1.8e10}{m^3}$, $Q_{\mathrm{out}} = \qty{0.3e10}{m^3}$ en
$k = 6$: voorbij de drempel van gips (3.8) maar niet die van haliet (10.6), dus
slaan calciumcarbonaat en gips neer. Met $Q_{\mathrm{in}} =
\qty{1.575e10}{m^3}$, $Q_{\mathrm{out}} = \qty{0.075e10}{m^3}$ en $k = 21$, voor
een ion dat in oplossing blijft, zoals magnesium: de pekel passeert de drempel van haliet,
dus haliet slaat ook neer, maar niet de magnesiumzouten (70).
(c)~Haliet vereist $Q_{\mathrm{in}}/(Q_{\mathrm{in}} - EA) \ge 10.6$, dat is
$Q_{\mathrm{in}} \le 10.6\,EA/9.6 = 1.104\,EA = \qty{1.66e10}{m^3}$ per jaar: de
instroom mag de verdamping hoogstens met ongeveer een tiende overtreffen.
\end{solution}

\begin{solution}{exo:b1:chemical-biogenic-sediments:11}
(a)~Zeewater bevat \qty{67.3}{mg} broomide per kilogram; bij 10.6~keer zo sterk
geconcentreerd bevat de pekel $67.3\times10.6 = \qty{713}{mg.kg^{-1}}$ en het eerste haliet
$0.032\times713 = \qty{22.8}{mg.kg^{-1}}$. (b)~Bij 30~keer, $67.3\times30 =
\qty{2019}{mg.kg^{-1}}$ in de pekel en \qty{64.6}{mg.kg^{-1}} in het haliet.
(c)~Het haliet neemt slechts 3.2~honderdsten van de concentratie van de pekel op, zodat
bijna al het broomide in oplossing blijft en geconcentreerd blijft terwijl het water
verdampt. (d)~De stijging registreert de concentratie van de pekel; de plotselinge daling
registreert haar verdunning, door een nieuwe instroom van zeewater in het bekken, of van zoet
water, voordat het haliet verder groeide.
\end{solution}

\begin{solution}{exo:b1:chemical-biogenic-sediments:12}
(a)~$t_c = \qty{32.65}{Myr}$, dus $0.02\times\num{32650} = \qty{653}{m}$ slijk
(\num{32650} intervallen van duizend jaar bij \qty{2}{cm}) en
$0.001\times\num{37350} = \qty{37}{m}$ rode klei. (b)~$t_c = (2500/350)^2 =
\qty{51.0}{Myr}$: \qty{1020}{m} slijk en \qty{19}{m} rode klei. Een dieper gelegen CCD
geeft een veel dikker slijk, en de rode klei van de diepe Atlantische vloer is dunner
dan die van de Stille Oceaan van dezelfde leeftijd. (c)~De CCD kan enige tijd dieper zijn geweest,
zodat de plaats er weer boven kwam; of de plaat bracht de plaats onder een gordel met hoge
productiviteit, zoals de evenaar, waar meer carbonaat neerregende; of de laag werd
meegevoerd door een afschuiving of een \omterm{def:b1:clastic-sediments:turbidite}{troebelingsstroom} van een hoger
deel van de zeebodem.
\end{solution}

\begin{solution}{pb:b1:chemical-biogenic-sediments:1}
\textbf{1.} $1030\times1000 = \qty{1.03e6}{kg}$ zeewater, dat
$1.03\times10^6\times0.035165 = \qty{36200}{kg}$ zout bevat.

\textbf{2.} Gips $1.03\times10^6\times1.581\times10^{-3} = \qty{1628}{kg}$;
haliet $1.03\times10^6\times0.02741 = \qty{28230}{kg}$.

\textbf{3.} Calciumcarbonaat $113.3/2710 = \qty{0.042}{m}$; gips $1628/2300 =
\qty{0.71}{m}$; haliet $28230/2170 = \qty{13.0}{m}$.

\textbf{4.} Het anhydriet heeft de massa $1628\times136.14/172.17 =
\qty{1288}{kg}$ en de dikte $1288/2970 = \qty{0.43}{m}$.

\textbf{5.} Droog is het zout $0.110 + 1.250 + 27.41 + 6.34 = \qty{35.11}{g}$ per
kilogram, waarvan haliet $27.41/35.11 = \qty{78}{\%}$ is: natrium en chloride
maken het meeste van het opgeloste zout uit, en haliet neemt er bijna alle op.

\textbf{6.} Van onder naar boven: calciumcarbonaat, gips, haliet, de bitterzouten.
Elk \omterm{def:b1:rock-forming-minerals:mineral}{mineraal} begint neer te slaan wanneer de pekel zijn verzadiging bereikt, het minst oplosbare
eerst, en legt zich op wat eerder is afgezet.

\textbf{7.} De diepte van de pekel is $1000/\mathrm{CF}$: \qty{556}{m} voor calciumcarbonaat,
\qty{263}{m} voor gips, \qty{94}{m} voor haliet.

\textbf{8.} $1 - 1/10.6 = \qty{90.6}{\%}$.

\textbf{9.} $1000/70 = \qty{14}{m}$ voor magnesiumsulfaat, $1000/90 = \qty{11}{m}$
voor de kaliumzouten.

\textbf{10.} Ongeveer \qty{1.10}{g.cm^{-3}} wanneer calciumcarbonaat vormt, 1.13 met
gips, 1.22 met haliet en meer dan 1.30 met de bitterzouten, tegenover 1.03 voor
zeewater. Binnenstromend zeewater drijft op de dichte pekel, die onderin blijft; in een
afgesloten bekken kan de pekel onder de instroom terugstromen naar de oceaan.

\textbf{11.} Hij is slechts \qty{4}{cm} dik, tegenover \qty{13}{m} haliet; en hij
vormt zich aan de randen, waar het binnenstromende water het eerst concentreert, meer dan
in het midden, dus is hij vermengd met de basis van het gips of erdoor vervangen.

\textbf{12.} Het haliet neemt slechts een klein deel van het broomide van de pekel op
(oefening~\ref{exo:b1:chemical-biogenic-sediments:11}), dus wordt de pekel
rijker aan broomide naarmate het haliet groeit, en neemt elke nieuwe laag haliet meer op.

\textbf{13.} $1500/13.01 = 115.3$: 116 vullingen.

\textbf{14.} $116\times0.434 = \qty{50}{m}$ anhydriet, niet als één bed maar als
116 lagen, elk iets minder dan een halve meter dik, aan de basis van het haliet van
hun vulling: een cyclische opeenvolging van anhydriet en haliet.

\textbf{15.} Per vulling $1.03\times10^6\times6.34\times10^{-3}/2000 =
\qty{3.3}{m}$, dus \qty{379}{m} voor 116 vullingen. Het zijn de meest oplosbare
zouten: de volgende instroom van zeewater lost ze weer op, en een bekken waarvan de pekel
voor het einde uitstroomt, bereikt ze nooit.

\textbf{16.} Bij elke bijvulling wordt de pekel verdund: het broomidegehalte daalt
scherp aan de basis van het nieuwe haliet en stijgt daarboven weer, een zaagtandprofiel
met één tand per vulling.

\textbf{17.} $2000/13.01 = 154$ vullingen.

\textbf{18.} $5000\times365.25\times86400 = \qty{1.578e11}{m^3}$ per jaar, verdeeld
over $\qty{1.0e11}{m^2}$: $E = \qty{1.58}{m}$ per jaar.

\textbf{19.} $1.578\times10^{11}\times1030\times0.035165 = \qty{5.72e12}{kg}$ zout
per jaar.

\textbf{20.} $1.578\times10^{11}\times1030\times0.02741/2170 =
\qty{2.05e9}{m^3}$ haliet per jaar, dat is \qty{0.0205}{m}, ongeveer \qty{2}{cm}
per jaar over de bodem.

\textbf{21.} $1500/0.0205 = \num{73000}$ jaar.

\textbf{22.} Het niveau daalt met \qty{1.58}{m} per jaar; \qty{1000}{m} water is
in $1000/1.578 = 634$ jaar verdwenen.

\textbf{23.} Haliet verschijnt wanneer \qty{90.6}{\%} van het water is verdwenen, na
$905.7/1.578 = 574$ jaar.

\textbf{24.} De \qty{1.5}{km} haliet kostte het zout van 116 vullingen van het
bekken (115.3, naar boven afgerond), of die nu na elkaar kwamen terwijl het bekken
droogliep en opnieuw vulde, of continu over de drempel, zoals in
vraag~21, tegen de snelheid van één volle zee elke
$1000/1.578 = 634$ jaar.
\end{solution}
